\section{Analysis of Recovered Dependency Graphs}\label{sec:findings}

In this section, we analyze the graph recovered by \system{} and highlight findings that would be difficult to surface from individual sources alone. One advantage of representing recursive dependencies as a structured graph is that audit questions become executable queries rather than manual document searches. We can issue graph/SQL-style queries over the recovered graph: find released models that transitively depend on generators with restrictive licenses; identify benchmarks that appear both as evaluation targets and as upstream training seeds, prompts, or concept sources; locate operations supported by only one source type; or detect stages where papers, cards, and code provide conflicting evidence. The findings below were surfaced through this query-driven audit workflow: structured queries identify candidate risk patterns, and the attached source anchors make those candidates verifiable.

Across the analyzed releases, \system{} recovers a large, heterogeneous, and deeply recursive dependency graph. Table~\ref{tab:graph-scale} reports the scale and depth of the recovered lineage for each target model, while Table~\ref{tab:graph-composition} summarizes the aggregate graph composition. The graph contains both model and dataset artifacts, reflecting that model dependencies often enter indirectly through generated, filtered, or rewritten datasets rather than only through base-model inheritance.

\begin{table}[t]
\centering
\small
\caption{Scale and depth of recovered dependency graphs. Ancestors count unique transitive upstream artifacts reachable from each target model.}
\label{tab:graph-scale}
\begin{tabular}{lcc}
\toprule
Target model & Ancestors & Max depth \\
\midrule
Olmo-3-Instruct & 618 & 8 \\
Olmo-3-Think & 675 & 11 \\
Olmo-3-Base & 601 & 11 \\
Nemotron-3-Super & 667 & 6 \\
Nemotron-3-Nano-Base & 362 & 5 \\
DR-Tulu & 33 & 4 \\
SmolLM3-Base & 76 & 2 \\
\midrule
\textbf{Aggregate graph} & \multicolumn{2}{c}{2,844 nodes, 14,769 edges, 51,456 anchors} \\
\bottomrule
\end{tabular}
\end{table}

\begin{table}[t]
\centering
\small
\caption{Aggregate composition of the recovered graph after deduplication. Artifact nodes are grouped into identity clusters, and evidence anchors record source excerpts supporting recovered relationships.}
\label{tab:graph-composition}
\begin{tabular}{lr}
\toprule
Statistic & Count \\
\midrule
Source lattices merged & 193 \\
Relation artifacts merged & 715 \\
Artifact identity clusters & 2,545 \\
Artifact nodes & 2,844 \\
\quad Dataset nodes & 1,472 \; (51.8\%) \\
\quad Model nodes & 1,070 \; (37.6\%) \\
\quad Off-lattice nodes & 299 \; (10.5\%) \\
\quad Entity nodes & 3 \\
\quad Concept / Misc.\ nodes & 0 \\
\midrule
Dependency edges & 14,769 \\
Evidence anchors & 51,456 \\
Flagged cross-source conflicts & 60,708 \\
\bottomrule
\end{tabular}
\end{table}

Table~\ref{tab:relation-types} shows that evaluation and training-data relationships dominate the recovered graph, while model-mediated data construction remains a substantial fraction of recovered edges. This supports the central motivation for recursive lineage tracing: modern models often depend on other models not only through weights, but through the data those models generate, filter, rewrite, judge, or label. This is especially audit-relevant given recent evidence that model-generated data can transmit model-specific behavioral traits even when explicit references are filtered out, including through code and reasoning traces~\cite{cloud2025subliminallearninglanguagemodels}.

\begin{table}[t]
\centering
\small
\caption{Distribution of recovered dependency relation types. Percentages are computed over 14,769 recovered edges. \sewon{Potentially separate direct vs. indirect?}\sewon{Actually a little confused by this table. Does the last row mean that it's model-model dependency? Then, does it mean everything else is model-data dependency? (I assume some of them is model-data, some others are data-model) The other thing I just realize is if Model A generates Data B, and Data B is used to train Model C, then it will only be counted as model-data and data-model dependency but not model-model dependency at least in this table, because it's not like there is no direct edge? I wonder if there is a better way...}}
\label{tab:relation-types}
\begin{tabular}{lrr}
\toprule
Relation type & Count & Share \\
\midrule
\texttt{used\_for\_evaluation} & 6,235 & 42.2\% \\
\texttt{trained\_on} & 4,772 & 32.3\% \\
\texttt{generated\_by} & 1,260 & 8.5\% \\
\texttt{inspired\_by} & 673 & 4.6\% \\
\texttt{trained\_from} & 554 & 3.8\% \\
\texttt{used\_for\_ablation} & 510 & 3.5\% \\
\texttt{filtered\_by} & 480 & 3.3\% \\
\texttt{transformed\_by} & 252 & 1.7\% \\
\texttt{decontaminated\_against} & 23 & 0.2\% \\
\texttt{quantized\_from} & 5 & $<0.1$\% \\
\texttt{embedded\_by} & 4 & $<0.1$\% \\
\texttt{merged\_from} & 1 & $<0.1$\% \\
\midrule
\textbf{Model-mediated generation/filtering/transformation} & \textbf{1,992} & \textbf{13.5\%} \\
\bottomrule
\end{tabular}
\end{table}

The recovered evidence is also highly fragmented across source types. As shown in Table~\ref{tab:source-type-distribution}, most operations are supported by only one source class, meaning that systems restricted to a single document type would miss a large fraction of recoverable lineage.

\begin{table}[t]
\centering
\small
\caption{Source-type distribution of recovered operations. Single-source operations are supported by evidence from only one source class; multi-source operations are supported by anchors from multiple source classes.}
\label{tab:source-type-distribution}
\begin{tabular}{lrr}
\toprule
Source support & Count & Share \\
\midrule
Only Hugging Face cards & 885 & 6.0\% \\
Only training code / scripts & 3,712 & 25.2\% \\
Only PDFs / technical reports & 4,817 & 32.8\% \\
Only release blogs / docs & 61 & 0.4\% \\
Only other docs & 983 & 6.7\% \\
Multiple source types & 4,243 & 28.9\% \\
\midrule
Total operations & 14,701 & 100\% \\
\bottomrule
\end{tabular}
\end{table}

\paragraph{Hidden external-model dependencies}

A core advantage of recursive dependency tracing is that it surfaces upstream models that never appear in the final model's own paper or card. These dependencies are often mediated by intermediate datasets, filters, classifiers, teachers, or tools, and therefore only become visible after following several evidence-backed hops across documents.

\begin{itemize}
    \item \textbf{SmolLM3 inherits a Llama-3-70B-mediated pretraining dependency through FineMath.}
    SmolLM3 documents FineMath as part of its pretraining mixture, and the FineMath card describes the use of a classifier to score and filter mathematical web data. Only by following the classifier card does the upstream dependency become visible: the \texttt{finemath-classifier} was trained on annotations generated by \texttt{Llama-3-70B-Instruct}. Thus, SmolLM3's pretraining pipeline is transitively shaped by Llama-3 through an intermediate filter, even though the SmolLM3 card does not name Llama-3 as a data source.

    \item \textbf{DR-Tulu inherits a Claude dependency through ScholarQA trajectories.}
    The DR-Tulu paper states that Ai2 ScholarQA trajectory data was used to create SFT data, but does not name the model used to generate those trajectories. Following the ScholarQA implementation reveals that its generation pipeline uses Claude Sonnet 3.7 by default. The graph therefore surfaces a transitive model-mediated dependency: \texttt{Claude Sonnet 3.7} $\rightarrow$ \texttt{Ai2 ScholarQA trajectories} $\rightarrow$ \texttt{DR-Tulu SFT}.

    \item \textbf{Nemotron-3-Content-Safety inherits DeepSeek-R1-generated content through NemoGuard.}
    The Nemotron-3-Content-Safety card describes the model as extending Llama-3.1-NemoGuard 8B Content Safety. The NemoGuard card identifies its base/fine-tuning lineage, while the NaturalReasoning documentation indicates that a subset of training questions was answered using DeepSeek-R1. Combining these sources reveals a transitive DeepSeek-R1-generated-data dependency that is not visible from the Nemotron-3-Content-Safety card alone.

    \item \textbf{OLMo-3 RL-Zero variants inherit Qwen2.5-Coder through intermediate checkpoints and data construction.}
    OLMo-3 RL-Zero models are trained on Dolci RL-Zero mixtures, which derive from earlier OLMo-3 checkpoints and code-oriented data construction. Tracing those intermediate artifacts back to the OLMo-3 midtraining/data-generation pipeline surfaces Qwen2.5-Coder-32B-Instruct as a transformation model used upstream. The dependency is not stated in the downstream RL-Zero cards, but emerges by following checkpoint and dataset lineage.

\end{itemize}

These examples show why final-model documentation is insufficient for model-mediated dependencies. A model card can accurately describe the immediate dataset or teacher while omitting the upstream generator, judge, or filter that shaped that artifact several steps earlier.

\paragraph{Train/eval coupling.}

The graph also surfaces relationships between training data and evaluation benchmarks that are not necessarily direct test-set leakage. Instead, these are cases of structural train/eval coupling: benchmark training splits, auxiliary data, prompts, validation splits, or concepts are transformed into training data while the same benchmark family remains an evaluation target. Such relationships can evade surface-form decontamination because the relevant connection is lineage-based rather than text-overlap-based.

\begin{itemize}
    \item \textbf{IFEval prompts are used in OLMo-3 training and IFEval remains an evaluation target.}
    OLMo-3 uses IF-RLVR prompts sampled from IFEval and IFBench-Train, while IFEval is also used in evaluation. This creates a direct audit question about whether evaluation prompts were excluded from training prompts.

    \item \textbf{SWE-bench and SWE-rebench appear in both training and evaluation roles within the Nemotron family.}
    Nemotron training artifacts include SWE-bench/SWE-rebench-derived resources, while Nemotron model cards or evaluation code report evaluations on SWE-bench-style benchmarks. This is stronger than cross-model reuse because the train/eval roles occur within the same product family.

    \item \textbf{LiveCodeBench validation appears in training artifacts while LiveCodeBench remains an evaluation benchmark.}
    Nemotron training artifacts include a \texttt{livecodebench\_v5\_validation} resource. When paired with LiveCodeBench evaluation use in related releases, this surfaces a benchmark-split audit question that would be hard to see without aligning training YAMLs and eval configs.
\end{itemize}

These cases do not by themselves establish benchmark contamination. Rather, they show that evaluation independence is a graph property: to assess it, one must trace whether benchmark artifacts appear upstream as seeds, prompts, validation splits, concept sources, or synthetic-data inputs.

\paragraph{License obligations propagate through synthetic-data pipelines.}

License risk is not always visible from the license attached to a final dataset or model. Prior work argues that dataset legal risk cannot be assessed from license terms alone and instead requires tracing redistribution and lifecycle history~\cite{kim2025trustlicensesseedataset}; supply-chain audits similarly find that license labels can function as weak metadata signals when the underlying compliance payload is missing or fails to propagate~\cite{jewitt2026permissivewashingopenaisupply}. Our recovered graphs surface analogous risks in LLM pipelines, where obligations may originate from models used to generate, label, or filter a dataset rather than only from the final model card.

\begin{itemize}
    \item \textbf{SmolLM3's FineMath pipeline is transitively connected to Llama-3-70B.}
    SmolLM3 names FineMath as a pretraining source; FineMath describes a classifier-based filtering pipeline; and the classifier card reveals that the classifier was trained on Llama-3-70B-Instruct-generated annotations. This does not by itself establish a legal violation, but it surfaces a non-obvious license/provenance dependency in the pretraining pipeline.

    \item \textbf{OLMo-3's use of Llama-Nemotron Post-Training data creates a Llama-license audit surface.}
    The OLMo-3 paper lists Llama-Nemotron Post-Training data as part of the training mixture, while the NVIDIA dataset card provides more precise generator and license information. Only by joining the OLMo mixture description with the upstream dataset card does the Llama Community License surface for this slice of the training pipeline.

    \item \textbf{PrismMath contributes to Nemotron-3-Super through a Qwen-license-bearing dataset.}
    Nemotron-3-Super training artifacts reference PrismMath, and the PrismMath card warns that it may be subject to redistribution and use requirements in the Qwen License Agreement. The graph therefore surfaces a Qwen-related license-obligation path that is not visible from the final Nemotron model card alone.

    \item \textbf{Nemotron pretraining data contains stacked generator-license surfaces.}
    Nemotron pretraining datasets such as \texttt{Nemotron-Pretraining-Specialized-v1}, \texttt{Nemotron-CC-v2}, and \texttt{Nemotron-CC-Math-v1} explicitly name obligations from Qwen, DeepSeek, and Phi-4. The graph makes these obligations queryable as a stacked license/provenance surface inherited by downstream Nemotron models.
\end{itemize}

We treat these as license-obligation findings rather than legal determinations. \system{} identifies where obligations may propagate.

\paragraph{Hidden judge, filter, and synthesizer stacks shape training data.}

Many dependencies do not appear as explicit training datasets. Instead, upstream models act as judges, filters, classifiers, reward models, embedding models, or synthetic-data generators. These roles are often compressed in papers but described only in scripts, dataset cards, or manifests.

\begin{itemize}
    \item \textbf{OLMo-3 uses multiple Gemma-3 variants for pretraining filtering.}
    The OLMo-3 paper mentions Gemma-based document classification, but the data-pipeline scripts reveal multiple distinct Gemma variants across model sizes and instruction/pretraining forms. The graph therefore exposes the full filtering panel rather than a single family-level mention.

    \item \textbf{Qwen3-32B is a pervasive judge across OLMo-3 RL variants.}
    The OLMo-3 paper notes the use of Qwen3-32B as an LM judge, while training scripts show that this judge is used across multiple RL-Zero specializations and post-training variants. A graph query over judge-role edges reveals the pervasiveness of this dependency across the OLMo-3 family.

    \item \textbf{OLMo-3-Think uses multiple external filtering and judging stages.}
    OLMo-3-Think is shaped by several model-mediated filters or judges across different stages, including Qwen-family models and GPT-4.1. These dependencies are scattered across paper sections and training scripts, and no single document enumerates them as a combined filtering pipeline.

    \item \textbf{Nemotron-3-Super uses DeepSeek-V3.2 as a tool-call trajectory judge.}
    The main Nemotron-3 report names some synthetic-data and judge models, but the agentic dataset card identifies DeepSeek-V3.2 as a judge for filtering tool-call trajectories. The graph therefore surfaces an additional judge dependency hidden in dataset-level documentation.

    \item \textbf{Nemotron-3-Nano-Omni vision data is filtered by a multi-organization judge panel.}
    Vision-data documentation for Nemotron-3-Nano-Omni enumerates filtering/generation roles played by GPT-4o, gpt-oss-120b, Gemini, and multiple Qwen variants. This shows that multimodal releases can inherit even larger external-model panels than text-only releases, and that those panels may appear only in data-source documentation.
\end{itemize}

\paragraph{Released artifacts can differ from trained-on artifacts.}

Artifact names alone are not stable evidence of what was used in training. A graph with per-operation evidence makes it easier to identify when released artifacts differ from the training-time artifacts they are meant to represent.

\begin{itemize}
    \item \textbf{The OLMo-3-7B reproduction dataset was modified after training.}
    The dataset intended to reproduce OLMo-3-7B-1025 was modified after training: some olmOCR science PDFs were redacted and replaced with \texttt{[REMOVED]}. The dataset card warns that this affects reproducibility, but the model paper does not foreground the issue.

    \item \textbf{Nemotron-3-Super RL-Training-Blends uses a placeholder-resolution pattern.}
    Nemotron-Super-RL-Training-Blends includes placeholder rows for DAPO-Math and Skywork-OR1 contributions, requiring a separate data-preparation script to fetch and restore upstream data before the blend is usable. Joining the dataset artifact, data-prep documentation, and YAML config reveals that the released blend depends on mutable upstream sources.

    \item \textbf{The placeholder-resolution pattern recurs across Nemotron product lines.}
    Similar placeholder-resolution mechanisms appear in Nemotron-3-Nano and Nemotron-3-Nano-Omni data pipelines. Combining documentation across product lines suggests that this is not a one-off artifact issue, but a recurring release strategy in the Nemotron family.

    \item \textbf{OLMo-3's OpenThoughts2 data path records model-name filtering not foregrounded in the paper.}
    The OLMo-core midtraining YAML includes an OpenThoughts2-derived path containing \texttt{modelnamefilter}. This indicates that explicit model-name filtering was applied in the upstream data path, a detail not surfaced in the main OLMo-3 paper.
\end{itemize}

\paragraph{Documentation and code disagree on reproducibility-critical details.}

The graph also exposes inconsistencies across papers, model cards, dataset cards, YAML metadata, READMEs, and training scripts. These cases are often small individually, but they matter because they affect reproducibility and interpretation of the training pipeline.

\begin{itemize}
    \item \textbf{SmolLM3's pretraining mixture is much more granular than its card summary.}
    The SmolLM3 card summarizes pretraining with broad categories such as web, code, and math data. The YAMLs reveal many specific subsets, language-specific substitutions, and stage-specific mixture weights. The card is useful for high-level understanding, but the YAMLs are necessary for reproduction.

    \item \textbf{SmolLM3's MegaMath source includes a Qwen-generated subset hidden behind a coarse dataset name.}
    The card names MegaMath as a single source, while code-level configuration distinguishes web-derived MegaMath data from Qwen-generated QA pairs. This reveals an additional model-generated subset that is obscured by the coarse dataset label.

    \item \textbf{OLMo-3 release artifacts contain smaller but concrete metadata conflicts.}
    OLMo-3 artifacts disagree on details such as the number of checkpoints used for souping, model-card prose sometimes names the wrong sibling dataset, and per-source prompt counts differ across related dataset cards. These examples show that even highly open releases can contain local documentation inconsistencies.
\end{itemize}

Together, these findings show that recursive model-model dependency graphs are more than just a descriptive artifact; they provide a strong auditing layer and can surface key issues that were previously opaque due to the complexity of modern LLM development. Because \system{} uses reported information from public artifacts, these findings should be interpreted as a lower bound on the true dependency structure. Undocumented uses of private models, unreleased data mixtures, or internal filtering pipelines may introduce additional dependencies that are not recoverable from public evidence.
